\documentclass[10pt, a4paper]{article}
\usepackage{preamble}

\title{Algebra II}
\author{Luke Phillips}
\date{October 2025}

\begin{document}

\maketitle

\newpage

\tableofcontents

\newpage

\section{Rings,
subrings and fields}

\subsection{First examples
(non-examples)}

\begin{enumerate}[label = (\roman*)]
    \item $\Z = \{0, \pm 1, \pm 2, \dotsc\}$.

    \item $\Q = \{\frac{a}{b} \mid a, b \in \Z, b \neq 0\}$.

    \item $\R$,
    $\C$.
    In these cases we can divide by non-zero elements.

    \item Polynomials:
    Let $Q[x] = \{a_0 + a_1x + \dotsc + a_nx ^ n \mid a_i \in \Q\}$.

    \item Matrices:
    Let $n \in \N \coloneqq \{1, 2, 3, \dotsc\}$ and define $M_n(\Q) = \{n \times n\text{ matrices with entries in } \Q\}$.

    We can add/multiply matrices.

    \item Functions.
    Let $C(\R) = \{f : \R \to \R \mid f\text{ continuous}\}$.
    We can add,
    subtract and multiply functions:
    \[
    (f \pm g)(x) = f(x) \pm g(x),\qquad (f\cdot g)(x) = f(x)g(x).
    \]
\end{enumerate}

Informal definition of a ring:
A ring is a set of objects where we can add/subtract/multiply two objects whilst remaining inside the set.
If we can also divide,
then this structure is called a field.

All the examples in the above list are rings,
$\Q, \R, \C$ are fields.

Non-examples:
\begin{enumerate}[label = (\roman*)]
    \item $\N \coloneqq \{1, 2, 3, \dotsc\}$ is a not a ring,
    $1 - 2 = -1 \notin \N$.

    \item Let $V = \R ^ 3$.
    Is $V$ a ring under $+$ and cross product?

    No,
    as in general
    \[
    u \times (v \times w) \neq (u \times v) \times w.
    \]
    This missing property is called associativity.
\end{enumerate}

Summarising:
$\Z$ is a ring,
but not a field.
I.e. $\frac{1}{2} \notin \Z$.

$\Q[x]$ is a ring,
but not a field:
$\frac{1}{x} \notin \Q[x]$.

$M_n(\Q)$,
$C(\R)$ are rings.

$\N$,
$\R ^ 3, +, \times$ are not rings.

\subsection{Definition of a ring}

\textbf{Binary operations} are operations which take in two inputs such that if $R$ is a set,
a binary operation is a function $f : R \times R \to R$,
$(x, y) \mapsto f(x, y)$.
$+(x, y)$ is addition,
normally written as $x + y$.
$+$ should make $R$ into an abelian group
($x + y = y + x$).

For $R$ to be a ring,
we also need another binary operation $\cdot : R \times R \to R$,
written $x \cdot y$ or $xy$.

\begin{definition}[Ring]
    A ring $R$ is a set together with two binary operations $+, \cdot$,
    such that $+$ makes $R$ into an abelian group and such that
    
    \begin{enumerate}[label = (\roman*)]
        \item (Identity)
        \[
        \exists 1 \in R, 1 \cdot x = x \cdot 1 = x,\quad\forall x \in R.
        \]

        \item (Associativity)
        \[
        x(yz) = (xy)z,\quad\forall x, y, z \in R.
        \]

        \item (Distributivity)
        \[
        x(y + z) = xy + xz\text{ and } (y + z)x = yx + zx\qquad\forall x, y, z \in R.
        \]
    \end{enumerate}
\end{definition}

\textit{Note:
$R$ is an abelian group implies that for any $x \in R, -x \in R$,
$x - y \coloneqq x + (-y)$.}

\subsection{Integers modulo \texorpdfstring{$n$}{}}
Let $n \in \N$.
If we divide an integer by $n$ we always get a remainder among $0, 1, 2, \dotsc, n - 1$.
When $a \in \Z$ is thought of as a remainder modulo $n$,
we write it $\bar{a}$.
$\bar{a}$ cannot be an integer,
e.g. if $n = 2$,
$\bar{a} = \bar{0}, \bar{1}$ and $\bar{0} = \bar{2} = \dotsi$ but $0 \neq 2$ so $\bar{a}$ cannot be defined.
\[
\bar{0} = 0, \pm 2, \pm 4, \dotsi.
\]
\[
\bar{1} = 1, \pm 3, \pm 5, \dotsi.
\]
Define
\[
\bar{0} = \{x \in \Z\mid x \equiv 0 \pmod{2}\},\qquad\bar{1} = \{x \in \Z \mid x \equiv 1\pmod{2}\}.
\]
For general $n$,
we define
\[
\bar{a} = \{x \in \Z \mid x \equiv a \pmod{n}\}.
\]
This is called a residue class.
There are $n$ of these:
\[
\bar{0}, \bar{1}, \dotsc, \overline{n - 1}.
\]
We denote this set by $\Z / n$
(sometimes written $\Z_{n}$).

\begin{example}
    We will show that $\Z / n$ is a commutative ring
    (i.e. $xy = yx$).

    \begin{solution}
        We need two binary operations $+, \cdot$:
        \begin{align*}
            \bar{a} + \bar{b} &\coloneqq \overline{a + b}. \\
            \bar{a} \cdot \bar{b} \coloneqq \overline{ab}.
        \end{align*}
        Thus,
        to define $+$ and $\cdot$ we have chosen $a$ and $b$ in $\Z$,
        as "representing" $\bar{a}$ and $\bar{b}$.
        To be sure that these operations define a ring,
        we need to check that they are independent of the choice of representatives.

        We leave the task of checking $\Z / n$ is an abelian group
        (under $+$),
        since this will not be as important.

        Existence of a
        (multiplicative)
        identity.
        $\bar{1} \in \Z / n$.
        commutativity $(xy = yx)$,
        associativity $(x(yz) = (xy)z)$ and distributivity $(x(y + z) = xy + xz)$ all follow from the corresponding properties in $\Z$,
        e.g.,
        \[
        \bar{a}(\bar{b} + \bar{c}) = \bar{a}(\overline{b + c}) = \overline{a(b + c)} = \overline{ab + ac} = \overline{ab} + \overline{ac} = \bar{a}\bar{b} + \bar{a}\bar{c}.
        \]
    \end{solution}
\end{example}

\begin{example}
    Consider $\Z / 3 = \{\bar{0}, \bar{1}, \bar{2}\}$.
    Addition and multiplication tables:
    
    \begin{table}[H]
        \begin{tabular}{c|c|c|c}
            & $\bar{0}$ & $\bar{1}$ & $\bar{2}$ \\
            \hline
             $\bar{0}$ & $\bar{0}$ & $\bar{1}$ & $\bar{2}$ \\
             $\bar{1}$ & $\bar{1}$ & $\bar{2}$ & $\bar{0}$ \\
             $\bar{2}$ & $\bar{2}$ & $\bar{0}$ & $\bar{1}$
        \end{tabular}
    \end{table}
    
    \begin{table}[H]
        \begin{tabular}{c|c|c|c}
            & $\bar{0}$ & $\bar{1}$ & $\bar{2}$ \\
            \hline
             $\bar{0}$ & $\bar{0}$ & $\bar{0}$ & $\bar{0}$ \\
             $\bar{1}$ & $\bar{0}$ & $\bar{1}$ & $\bar{2}$ \\
             $\bar{2}$ & $\bar{0}$ & $\bar{2}$ & $\bar{1}$
        \end{tabular}
    \end{table}
\end{example}

\subsection{Subrings}

\begin{example}
    $\Z \subseteq \Q$,
    $\R \subseteq \C$,
    $\Q \subseteq \Q[x]$.
\end{example}

\begin{definition}[Subring]
    A subring $S$ of a ring $(R, +, \cdot)$ is a set $S \subseteq R$ such that
    
    \begin{enumerate}[label = (\roman*)]
        \item $0 \in S$,
        $1 \in S$.

        \item $a, b \in S \implies a + b \in S$.

        \item $a, b \in S \implies ab \in S$.

        \item $a \in S \implies -a \in S$.
    \end{enumerate}
\end{definition}

\begin{example}
    Show that $\Z[\sqrt{2}] = \{a + b\sqrt{2} \mid a, b \in \Z\}$ is a ring.

    \begin{solution}
        We can show that $\Z[\sqrt{2}]$ is a subring of $\R$.

        We can clearly see that $0, 1 \in Z[\sqrt{2}]$.
        To see associativity, distributivity and inverse.
        \begin{align*}
            (a + b\sqrt{2}) + (c + d\sqrt{2}) &= (a + b) + (b + d)\sqrt{2} \\
            (a + b\sqrt{2})(c + d\sqrt{2}) &= (ac + 2bd) + (ad + bc)\sqrt{2} \\
            -(a + b\sqrt{2}) &= -a - b\sqrt{2}.
        \end{align*}
        Thus,
        $\Z[\sqrt{2}]$ is a subring of $\R$.
        We do not have to check associativity or distributivity as these follow from the corresponding known properties in $\R$.
    \end{solution}
\end{example}

\begin{example}
    Consider the ring $\Z / 6$.
    It has $R = \{\bar{0}, \bar{2}, \bar{4}\}$ as a subset which satisfies the properties of being a subring
    (apart from $\bar{1} \in R$).
    It is a ring:
    $\bar{4}$ works as the multiplicative identity,
    e.g. $\bar{4} \cdot \bar{0} = \bar{0}$.
\end{example}





















\end{document}